\section{Strategy 4: the direct nine-point obstruction}
\label{sec:ab-core}

This section proves a center-class-independent theorem.  Its hypotheses are
that all six vertex roles are $\Vdzero$, that they cover the full boundary,
and that exactly one actual boundary row is supercritical.  The theorem will
therefore apply both to the $\CEzero$ branch and to the zero-gap subcases of
$\CEone$ and $\CEtwo$.  Six common radial points and three asymmetric points
are excluded directly from all six actual vertex roles and are therefore
forced into the center role.  Only the unique supercritical row uses an
$AB$-union; no nonsupercritical model core is introduced.

\subsection{Corner-clipped $AB$-unions}

Let $e_A,e_B$ be unit vectors with angle $120^\circ$, put
\[
 O=0,\qquad A=ae_A,\qquad B=be_B,
 \qquad
 C=\{ue_A+ve_B:u,v\ge0\},
\]
and define
\begin{equation}
 \mathcal R(a,b)=
 C\cap\bigcup\{T:T\text{ is a closed unit equilateral triangle and }
 O,A,B\in T\}.                               \label{eq:ab-union}
\end{equation}
The finite description in Theorem~\ref{thm:cert-caliper} applies to the
four-point set $\{O,A,B,x\}$ and makes membership in \eqref{eq:ab-union} an
exact finite semialgebraic condition of degree at most four.

The strict supercritical branch admits a simpler closed form.  In the next
statement write $x=ue_A+ve_B$ and
$\lVert x\rVert^2=u^2+v^2-uv$.

\begin{theorem}[Strict supercritical $AB$-union]
\label{thm:strict-ab-union}
Assume
\[
 a+b>1,\qquad \rho=a^2+ab+b^2<1,
\]
put $h=\sqrt3/2$, $D=\sqrt{4\rho-3}$, and define
\begin{align*}
 \alpha&=\frac{h(a+2b-aD)}{2\rho},&
 \beta&=\frac{h(a-b+(a+b)D)}{2\rho},\\
 \gamma&=\frac{h(-a+b+(a+b)D)}{2\rho},&
 \delta&=\frac{h(2a+b-bD)}{2\rho}.
\end{align*}
Then all four coefficients are positive and
\begin{equation}
 \mathcal R(a,b)=
 (C\cap\mathcal D_A\cap\mathcal H_A)
 \cup(C\cap\mathcal D_B\cap\mathcal H_B),       \label{eq:strict-ab-identity}
\end{equation}
where
\begin{align*}
 \mathcal D_A&=\{(u,v):(u-a)^2+v^2-(u-a)v\le1\},\\
 \mathcal D_B&=\{(u,v):u^2+(v-b)^2-u(v-b)\le1\},\\
 \mathcal H_A&=\{(u,v):\alpha(u-a)+\beta v\le0\},\\
 \mathcal H_B&=\{(u,v):\gamma u+\delta(v-b)\le0\}.
\end{align*}
Its non-axis frontier consists, in counterclockwise order, of the
$A$-circle, the $A$-line, the $B$-line, and the $B$-circle.
\end{theorem}

\begin{proof}
The identities
\[
 4\rho-3(a+b)^2=(a-b)^2,
 \qquad
 (a+b)^2D^2-(a-b)^2=4\rho((a+b)^2-1)
\]
give $0<D<1$ and positivity of the four coefficients.

We apply the caliper criterion of Theorem~\ref{thm:cert-caliper}.  For a
point $x=(u,v)$ in the relative interior of $C$ but outside
$\triangle OAB$, the four hull vertices occur in the order $O,B,x,A$.
The two cone-axis hull edges cannot certify a unit triangle: their support
sums, divided by $h$, are respectively
\[
 \max\{a,u\}+\max\{b,v-u\},\qquad
 \max\{b,v\}+\max\{a,u-v\},
\]
and are at least $a+b>1$.  It remains to test the $Ax$ and $Bx$ calipers.

For the $Ax$ caliper put $p=u-a$ and
$r^2=p^2+v^2-pv$.  Exact evaluation of the three supports gives
\begin{equation}
 \frac{G_A}{h}=
 \frac{\max\{r^2,-ap+(a+b)v\}+N_A}{r},
 \quad
 N_A=\begin{cases}
 0,&p\le0,\\
 ap,&0\le p\le v,\\
 (a+b)p-bv,&p\ge v.
 \end{cases}                                      \label{eq:Ax-caliper}
\end{equation}
No support label is omitted in these three sectors.  If $p\le0$,
$G_A\le h$ is equivalent to
\[
 r\le1,\qquad -ap+(a+b)v\le r.
\]
The second inequality is the half-plane condition in
\eqref{eq:strict-ab-identity}, because
\[
 r^2-(-ap+(a+b)v)^2
 =\frac{(cp+dv)(c'p+d'v)}{4\rho},
\]
where
\[
 c=a+2b-aD,\quad d=a-b+(a+b)D,
 \quad c'=a+2b+aD,\quad d'=a-b-(a+b)D,
\]
and $c,c',d>0>d'$.  Thus the low sector is exactly
$\mathcal D_A\cap\mathcal H_A$.  In the middle sector $0<p\le v$ one has
$r\le v$, whereas the numerator in \eqref{eq:Ax-caliper} is at least
$(a+b)v>r$, so no fit is possible.

In the high sector $p\ge v$, the inequalities $G_A\le h$ imply
\begin{equation}
 r^2+(a+b)p-bv\le r,\qquad bp+av\le r.          \label{eq:Ax-high}
\end{equation}
Write $x-A=ry(\theta)$, $0\le\theta\le60^\circ$, with cone coordinates
\[
 y(\theta)=te_A+we_B,\quad
 t=\frac2{\sqrt3}\cos(\theta-30^\circ),\quad
 w=\frac2{\sqrt3}\sin\theta.
\]
Then \eqref{eq:Ax-high} gives
\[
 r\le f(\theta)=1-(a+b)t+bw,\qquad S(\theta)=bt+aw\le1.
\]
Let $n_B$ be the unit vector with dual coordinates $(\gamma,\delta)$.
The identities
\[
 b\gamma+(a+b)\delta=h,\quad
 (a+b)\gamma+a\delta=\frac h2(1+D),\quad
 b\delta-a\gamma=\frac h2(1-D)
\]
show that the unique angle $\theta_B$ of this normal satisfies
$S(\theta_B)=1$ and $f(\theta_B)=(1-D)/2$.  The function $S$ has at most
one critical point, a maximum, so $S(\theta)\le1$ implies
$\theta\le\theta_B$.  On the feasible part $f,f'\ge0$; hence
$f(\theta)\cos(\theta_B-\theta)$ is increasing.  Consequently
\[
 \langle n_B,x-B\rangle
 \le a\gamma-b\delta+\frac h2(1-D)=0.
\]
The diameter bound also gives $\lVert x-B\rVert\le1$.  Thus every high
sector fit belongs to $\mathcal D_B\cap\mathcal H_B$.  Reflection exchanges
$a,u$ with $b,v$ and supplies the corresponding result for the $Bx$
caliper.  This proves \eqref{eq:strict-ab-identity} in the cone interior;
the axes follow by taking limits, since both sides are closed.

For completeness, the line--circle junctions are obtained by solving the
two displayed line and circle equations.  The line normals satisfy
\[
 \alpha^2+\alpha\beta+\beta^2=h^2,\qquad
 \gamma^2+\gamma\delta+\delta^2=h^2,
\]
and the line determinant is
\[
 \alpha\delta-\gamma\beta
 =\frac{3(1-D)(D+3)}{16\rho}>0.
\]
The successive cone slopes are
\[
 +\infty,\quad \frac2{1-D},\quad
 \frac{a(1-D)+2b}{2a+b(1-D)},\quad
 \frac{1-D}{2},\quad0;
\]
the two middle comparisons reduce to
$a(4-(1-D)^2)>0$ and $b(4-(1-D)^2)>0$.  This proves the asserted four-piece
order.
\end{proof}

\subsection{Exact-one handoffs and common demands}

\begin{lemma}[Exact-one handoff chain]
\label{lem:ab-extreme-jump}
Assume the six vertex roles cover $\partial H$ and exactly one actual row is
supercritical.  Rotate its index to $4$.  The strict handoffs supplied by
Theorem~\ref{thm:strict-handoff} satisfy
\[
 x_3<x_2<x_1<x_0<x_5<x_4.                       \label{eq:extreme-chain}
\]
With
\[
 a=1-x_3,\qquad b=x_4,\qquad p=1-b,\qquad q=1-a,
\]
one has
\begin{equation}
 0<a,b<1,\qquad a+b>1,\qquad a^2+ab+b^2<1,       \label{eq:ab-strict-domain}
\end{equation}
and every selected row obeys
\[
 a_i\ge p,\qquad b_i\ge q.
\]
\end{lemma}

\begin{proof}
Every selected row other than $4$ is strictly nonsupercritical, so
$x_i<x_{i-1}$ for $i\ne4$, which is precisely
\eqref{eq:extreme-chain}.  In particular $x_3$ is the global minimum and
$x_4$ the global maximum.  The two anchors of row $4$ lie in its open role;
their mutual squared distance is $a^2+ab+b^2$, strictly below the diameter
squared of the closed unit triangle.  This proves \eqref{eq:ab-strict-domain}.
Finally every $x_j\in[x_3,x_4]$, so
$1-x_{i-1}\ge1-x_4=p$ and $x_i\ge x_3=q$.
\end{proof}

\subsection{Six directly forced radial points}

Put $h=\sqrt3/2$ and
\[
 s=p+q,\qquad m=\min(p,q),\qquad M=\max(p,q),
 \qquad \chi=s^4-s^2+pq.
\]
The strict domain \eqref{eq:ab-strict-domain} gives
\begin{equation}
 pq>(1-s)(2-s),\qquad m>1-s,\qquad s>1-m.         \label{eq:pq-domain}
\end{equation}
Let $c_*=c_{\max}(p,q)$ be the exact radial envelope of
Proposition~\ref{prop:exact-admissible-set}.  On the present subcritical
domain, equation \eqref{eq:radial-envelope} becomes
\begin{equation}
 c_*=\begin{cases}
 C_L(m),&\chi\le0,\\[1mm]
 \displaystyle\frac{2M}{1+\sqrt{4s^2-3}},&\chi>0,
 \end{cases}                                    \label{eq:core-cstar}
\end{equation}
where $C_L(m)$ is the unique root in $[\sqrt3/2,1]$ of
\[
 z^4-z^2+mz-m^2=0.
\]
At $\chi=0$ the two expressions agree and equal $s$.

\begin{lemma}[Direct radial forcing]
\label{lem:symmetric-core-witness}
One has
\begin{equation}
 0<c_*<1,\qquad c_*>1-m.                         \label{eq:cstar-bounds}
\end{equation}
Assume that every $U_i$ is $\Vdzero$ and that
$U_C,U_0,\ldots,U_5$ cover $H$.  Then, for every $i$,
\[
 D_i=(1-c_*)V_i\in U_C.
\]
Consequently $U_C$ contains the centered closed disk
\begin{equation}
 \mathcal D_\eta=\{x:\lVert x\rVert\le\eta\},
 \qquad \eta=h(1-c_*).                           \label{eq:forced-disk}
\end{equation}
\end{lemma}

\begin{proof}
On the $C_L$ branch, $C_L(m)<1$ because the defining polynomial is positive
at $1$; if $s<\sqrt3/2$, then $c_*\ge\sqrt3/2>s>1-m$, while if
$s\ge\sqrt3/2$ its value at $s$ is $\chi\le0$, so $c_*\ge s>1-m$.
On the other branch put $E=\sqrt{4s^2-3}$.  The selector gives
$sE>M-m$, whence $c_*<s<1$.  Moreover, with
\[
 A_0=\frac{2M}{1-m}-1,
\]
one has
\[
 A_0^2-E^2=
 \frac{4(1-s)\{(1-m)(1-2m)+m(2-m)(1-s)\}}{(1-m)^2}>0.
\]
Thus $A_0>E$ and $c_*>1-m$.  This proves
\eqref{eq:cstar-bounds}.

Fix $i$.  If $D_i\in U_i$, openness supplies $\varepsilon>0$ with
$c_*+\varepsilon<1$ such that
$(1-c_*-\varepsilon)V_i\in U_i$.  The same role contains its two selected
handoff anchors.  Since their incoming and outgoing demands dominate $p$ and
$q$, convexity and passage to the closure give a closed unit triangle
realizing $(p,q,c_*+\varepsilon)$.  This contradicts the definition of
$c_*=c_{\max}(p,q)$.  Hence $D_i\notin U_i$.

If $D_i$ belonged to $U_{i-1}$ or $U_{i+1}$, a small plane ball inside that
open role would meet the adjacent radial arm $r_i$ in a positive-length
interval.  This contradicts the $\Vdzero$ definition.  For the remaining
roles put $r=1-c_*>0$.  Directly,
\[
 \lVert rV_i-V_{i\pm2}\rVert^2=1+r+r^2>1,
 \qquad
 \lVert rV_i-V_{i+3}\rVert=1+r>1.
\]
Since $V_j\in U_j$ and two points of an open unit triangle are at distance
strictly below $1$, these inequalities exclude $D_i$ from the other three
vertex roles.  Thus no vertex role contains $D_i$.  The point lies in
$\interior(H)$ by \eqref{eq:cstar-bounds}, so the assumed cover forces
$D_i\in U_C$.

The six points $D_i$ form a regular hexagon of circumradius $1-c_*$.  Their
convex hull contains its centered incircle of radius $h(1-c_*)$, and convexity
of $U_C$ gives \eqref{eq:forced-disk}.
\end{proof}

We next define the three asymmetric witnesses.  At $V_4$ use local
coordinates
\[
 P=V_4+u(V_5-V_4)+v(V_3-V_4),\qquad
 \mathsf q(u,v)=u^2+v^2-uv.
\]
The outgoing demand is $b$ and the incoming demand is $a$.  Accordingly,
set
\begin{align*}
 \alpha&=\frac{h(b+2a-bD)}{2\rho},&
 \beta&=\frac{h(b-a+(a+b)D)}{2\rho},\\
 \gamma&=\frac{h(-b+a+(a+b)D)}{2\rho},&
 \delta&=\frac{h(2b+a-aD)}{2\rho},
\end{align*}
where now $\rho=a^2+ab+b^2$ and $D=\sqrt{4\rho-3}$.  The two line pieces of
Theorem~\ref{thm:strict-ab-union} are
\[
 L_-(\lambda)=(b-\beta\lambda,\alpha\lambda),\qquad
 L_+(\mu)=(\delta\mu,a-\gamma\mu).
\]
They meet at $\lambda_* =\mu_*=8h\rho/(3(D+3))$ in
\begin{equation}
 J=\left(\frac{2b+a-aD}{D+3},
          \frac{b+2a-bD}{D+3}\right).           \label{eq:line-junction}
\end{equation}
Put
\[
 E_-=(1-b,2-b),\qquad E_+=(2-a,1-a).
\]
The actual adjacent handoffs have local centers
\[
 C_5(t)=(1+t,t),\qquad 1-a\le t\le b,
 \quad\text{and}\quad
 C_3(y)=(y,1+y),\qquad 1-b\le y\le a.
\]
The restrictions of the two unit-circle equations are
\begin{align}
 g_-(\lambda)&=\frac34\lambda^2
 -3(\alpha+b\beta)\lambda+3b^2-3b+2,\notag\\
 g_+(\mu)&=\frac34\mu^2
 -3(\delta+a\gamma)\mu+3a^2-3a+2.              \label{eq:witness-quadratics}
\end{align}
Let $\lambda_-$ and $\mu_+$ be their first roots; explicitly,
\begin{align*}
 \lambda_-&=2(\alpha+b\beta)
 -\frac23\sqrt{9(\alpha+b\beta)^2-9b^2+9b-6},\\
 \mu_+&=2(\delta+a\gamma)
 -\frac23\sqrt{9(\delta+a\gamma)^2-9a^2+9a-6}.
\end{align*}

\begin{lemma}[Moving adjacent-role signs]
\label{lem:fixed-line-signs}
Throughout \eqref{eq:ab-strict-domain},
\begin{equation}
 \lambda_*<\lambda_-<h^{-1},\qquad
 \mu_*<\mu_+<h^{-1}.                            \label{eq:first-root-order}
\end{equation}
Moreover $J_u,J_v<1/2$.  The point $L_+(\mu_+)$ satisfies
\[
 (L_+(\mu_+))_u+(L_+(\mu_+))_v<3-2a,
\]
and the reflected bound holds for $L_-(\lambda_-)$.  Writing
$F(P,t)=\mathsf q(P-C_5(t))-1$, one has, for every $t\in[1-a,b]$,
\begin{equation}
 F(L_+(\mu_+),t)\ge0,\qquad F(J,t)>0,\qquad
 F(L_-(\lambda_-),t)>0.                           \label{eq:moving-right-signs}
\end{equation}
Under the reflection
$(u,v,a,b)\leftrightarrow(v,u,b,a)$, the analogous signs hold for every
$y\in[1-b,a]$: the first-root point $L_-(\lambda_-)$ has nonnegative sign
against $C_3(y)$, and the other two witnesses have strictly positive sign.
\end{lemma}

\begin{proof}
Put
\[
 U=a+b-1,\qquad z=a-b,\qquad
 D^2=z^2+6U+3U^2.
\]
Then $0<U<1/6$ and $z^2<\Omega:=1-6U-3U^2$.  We give the
fixed-line calculation for the circle centered at
$E_+=(2-a,1-a)$; reflection gives the calculation for $E_-$.  Put
\[
 F(P,t)=\mathsf q(P-(1+t,t))-1,\qquad t_0=1-a,
\]
so this circle is $F(P,t_0)=0$.  At the junction, exact expansion gives
\[
 (D+3)^2F(J,t_0)=3D(z^2+Uz+1-3U)+P_J,
\]
where
\[
 P_J=z^4+5z^2+6Uz^2+3U^2z^2+9Uz+3U+15U^2>0;
\]
indeed $z^2+Uz+1-3U\ge1-3U-U^2/4>0$, and
$5z^2+9Uz+15U^2$ is positive definite.  Thus $F(J,t_0)>0$.
Also
\[
 J_u+J_v=\frac{(1+U)(3-D)}{3+D}<1.
\]
Indeed this inequality is $3U<D(2+U)$, which follows from
$D^2\ge3U(2+U)$ and $(2+U)^3>3U$.  Since
\[
 \frac{\partial F}{\partial t}(P,t)=1+2t-P_u-P_v,
\]
$F(J,t)$ is strictly increasing for $t\ge0$, and therefore is positive
throughout $[1-a,b]$.

On the opposite line $L_-$, the derivative at the junction, after
multiplication by a positive factor, has at $t=1-a$ the numerator
\[
 N_0=D(9+3z^2+6Uz-9U^2)+P_0,
\]
where
\[
\begin{aligned}
 P_0={}&18U^3+6U^2z^2+63U^2+18Uz^2+18Uz+54U\\
      &+2z^4+9z^2+9.
\end{aligned}
\]
The coefficient of $D$ is at least $9-6U-9U^2>0$, while
$P_0\ge9+54U-18U>0$.  Hence $N_0>0$.  At the other endpoint $t=b$, the
corresponding numerator is
\[
 N_1=D(9+3z^2+6U-3U^2)+P_1,
\]
where
\[
\begin{aligned}
 P_1={}&9U^4-3U^3z+45U^3+9U^2z^2-6U^2z+72U^2\\
      &-Uz^3+21Uz^2+9Uz+45U+2z^4+9z^2+9.
\end{aligned}
\]
Its $D$-coefficient is positive and, using $|z|<1$,
\[
 P_1\ge9+45U-9U-U-6U^2-3U^3>0.
\]
The junction derivative is affine in $t$, so it is positive throughout
$[1-a,b]$.  Since the restriction of $F$ to $L_-$ is a convex quadratic,
is positive at the junction, and is increasing there, it has no intersection
on the opposite line side for any actual adjacent handoff.

On the correct line $L_+$, the value at its outer endpoint $h^{-1}$ has,
after multiplication by a positive factor, numerator
\[
 A_0=P_2-4DU(1+U+z),
\]
where
\[
\begin{aligned}
 P_2={}&3U^4+6U^3z+6U^3+4U^2z^2+12U^2z+12U^2\\
      &+2Uz^3+6Uz^2+2Uz+6U+z^4+2z^2-3.
\end{aligned}
\]
For fixed $U$, replace the odd powers of $z$ by $x=|z|$.  The resulting
polynomial $P_2^+(U,x)$ satisfies
\[
 \partial_xP_2^+
 =2(3U^3+4U^2x+6U^2+3Ux^2+6Ux+U+2x^3+2x)>0,
\]
and therefore its supremum for $z^2<\Omega$ occurs at $x=\sqrt\Omega$.
There
\[
 P_2^+(U,\sqrt\Omega)=4U(U+\sqrt\Omega-3)<0.
\]
As $1+U+z=2a>0$, this gives $A_0<0$.  The positive junction value and
negative endpoint value force the first root strictly between them, proving
the $L_+$ half of \eqref{eq:first-root-order}; reflection proves the other
half.

For the coordinate bounds, the exact identities
\[
 D^2-(3U-z)^2=6U(1+z-U),\qquad
 D^2-(3U+z)^2=6U(1-z-U)
\]
give $J_u,J_v<1/2$.  It remains to control the coordinate sum along $L_+$.
For $E=L_+(h^{-1})$, the sign of $3-2a-E_u-E_v$ is the sign of
\[
 D(6U+2z+6)+B(U,z),
\]
where
\[
\begin{aligned}
 B(U,z)={}&-9U^3-9U^2z-9U^2-3Uz^2-18Uz+3U\\
          &-3z^3+3z^2-3z+3.
\end{aligned}
\]
Here
\[
 B_z=-3(3z^2+2Uz+6U-2z+3U^2+1)<0,
\]
because the quadratic in parentheses has discriminant
$-32U^2-80U-8<0$.  Hence its minimum on $z^2<\Omega$ is at
$z=\sqrt\Omega$, where
$B(U,\sqrt\Omega)=6(1-3U-\sqrt\Omega)>0$ because
$(1-3U)^2-\Omega=12U^2$.  Thus $E_u+E_v<3-2a$.
Also $J_u+J_v<1<3-2a$, and the coordinate sum is affine on the
line segment, so the same strict bound holds at the first root.  Now
\[
 \frac{\partial F}{\partial t}(P,t)=1+2t-P_u-P_v.
\]
Since $F(L_+(\mu_+),1-a)=0$, the coordinate-sum bound makes its $t$-derivative
strictly positive for $t\ge1-a$.  Together with the no-opposite-line result,
this proves \eqref{eq:moving-right-signs}.  Reflection gives both the bound on
$L_-$ and the complete $C_3(y)$ statement.
\end{proof}

Define the local and global witnesses by
\begin{equation}
 Q_-^{\rm loc}=L_-(\lambda_-),\qquad
 Q_0^{\rm loc}=J,\qquad
 Q_+^{\rm loc}=L_+(\mu_+),                       \label{eq:asym-witnesses}
\end{equation}
and by the displayed conversion from $(u,v)$ to the plane.

\begin{lemma}[Direct asymmetric forcing]
\label{lem:asymmetric-core-witness}
Assume that $U_C,U_0,\ldots,U_5$ cover $H$.  Then
\[
 Q_-,Q_0,Q_+
 \in\interior(H)\setminus\bigcup_{i=0}^5U_i,
\]
and consequently all three witnesses belong to $U_C$.
\end{lemma}

\begin{proof}
The root order and positivity of the line coordinates give
$0<u,v<1$ for all three points.  Each belongs to $\mathcal R(b,a)$ and hence
lies in a closed unit triangle containing $V_4$, so
$\mathsf q(u,v)\le1$; hence
$|u-v|<1$ and all three points lie in $\interior(H)$.

The closure of $U_4$ is one of the triangles represented by
$\mathcal R(b,a)$, since it contains $V_4,X_3,X_4$.  All three witnesses lie
on its exact non-axis frontier.  If one belonged to the open triangle $U_4$,
a sufficiently small plane ball about it would lie both in $U_4$ and in the
interior of the corner cone, making the witness an interior point of
$\mathcal R(b,a)$.  This contradiction excludes $U_4$.

For the roles $U_0,U_1,U_2$, their contained vertices have local coordinates
\[
 \begin{array}{c|ccc}
 i&0&1&2\\ \hline
 V_i&(2,1)&(2,2)&(1,2).
 \end{array}
\]
For $P=(u,v)$,
\[
 \lVert P-(2,1)\rVert^2-1=\mathsf q(u,v)+2-3u,
 \quad
 \lVert P-(1,2)\rVert^2-1=\mathsf q(u,v)+2-3v.
\]
The inequalities $J_u,J_v<1/2$ and the line order give distance greater than
one from $V_0$ for $Q_-,Q_0$ and from $V_2$ for $Q_0,Q_+$.  For the remaining
pair, the first-root circle and the sum bound in
Lemma~\ref{lem:fixed-line-signs} give
\[
 \lVert Q_+-V_0\rVert^2-1
 =a(3-a-(Q_+)_u-(Q_+)_v)>a^2,
\]
and, by reflection, $\lVert Q_--V_2\rVert^2-1>b^2$.  All three points are
farther than one from $V_1=(2,2)$, because writing $u=1-\xi$, $v=1-\eta$
gives
\[
 \lVert P-V_1\rVert^2
 =1+\xi+\eta+\xi^2+\eta^2-\xi\eta>1.
\]
Since two points in one open unit equilateral triangle have distance strictly
less than one, these inequalities exclude $U_0,U_1,U_2$.

The actual handoff $X_5=C_5(x_5)$ belongs to $U_5$, and
$1-a<x_5<b$ by \eqref{eq:extreme-chain}.  The three signs in
\eqref{eq:moving-right-signs} therefore put every witness at distance at
least one from $X_5$, excluding $U_5$.  Similarly, if $y=1-x_2$, then
$1-b<y<a$, $X_2=C_3(y)\in U_3$, and the reflected signs exclude $U_3$.
Thus no vertex role contains any witness.  The assumed cover now forces all
three interior witnesses into $U_C$.
\end{proof}

The nine directly forced points are
$D_0,\ldots,D_5,Q_-,Q_0,Q_+$.  Under a putative cover,
Lemmas~\ref{lem:symmetric-core-witness} and
\ref{lem:asymmetric-core-witness} therefore force the compact set
\begin{equation}
 K_{\rm wit}(a,b)=
 \mathcal D_\eta\cup\{Q_-,Q_0,Q_+\}
 \subset U_C.                                    \label{eq:forced-witness-set}
\end{equation}

\subsection{The enclosure inequality}

For a compact set $K$, let $\Lambda(K)$ be the least side length of a closed
equilateral triangle containing $K$.  With $R$ denoting rotation through
$2\pi/3$,
\begin{equation}
 \Lambda(K)=\frac1h\min_{\lVert n\rVert=1}
 \sum_{j=0}^2h_K(R^jn).                          \label{eq:lambda-support}
\end{equation}

\begin{theorem}[Witness enclosure]
\label{thm:witness-enclosure}
For every $(a,b)$ in \eqref{eq:ab-strict-domain},
\[
 \Lambda(K_{\rm wit}(a,b))\ge1.
\]
\end{theorem}

\begin{proof}
The disk alone gives
$\Lambda(K_{\rm wit})\ge3\eta/h=3(1-c_*)$, so the assertion is automatic
when $c_*\le2/3$.  Assume henceforth $c_*>2/3$.

We first remove the square radicals in the first-root witnesses.  Normalize
$\bar\alpha=\alpha/h$, and similarly for the other three coefficients, and
put
\[
 \xi=\lambda/h,\quad \zeta=\mu/h,\quad
 \xi_*=\zeta_*=\frac{8\rho}{3(D+3)}.
\]
Let $g_-,g_+$ be the rescaled forms of
\eqref{eq:witness-quadratics}, and take one Newton step at the junction:
\[
 \widehat\xi_-=\xi_*-\frac{g_-(\xi_*)}{g_-'(\xi_*)},\qquad
 \widehat\zeta_+=\zeta_*-\frac{g_+(\zeta_*)}{g_+'(\zeta_*)}.
\]
Since each quadratic is strictly convex, positive and decreasing at the
junction, its tangent zero lies strictly between the junction and the first
root.  Define, in local coordinates,
\begin{align*}
 A&=\left(b-\frac34\bar\beta\widehat\xi_-,
          \frac34\bar\alpha\widehat\xi_-\right),\\
 B&=J,\\
 C&=\left(\frac34\bar\delta\widehat\zeta_+,
          a-\frac34\bar\gamma\widehat\zeta_+\right),
\end{align*}
and convert them to global vectors.  Then
\[
 A\in(Q_0,Q_-),\qquad C\in(Q_0,Q_+).
\]
Consequently, for
\[
 \widehat K=\conv(\mathcal D_\eta\cup\{A,B,C\}),
\]
we have $\widehat K\subseteq\conv(K_{\rm wit})$ and
\begin{equation}
 \Lambda(K_{\rm wit})\ge\Lambda(\widehat K).    \label{eq:tangent-reduction}
\end{equation}

Put $S=\widehat K+R\widehat K+R^2\widehat K$.  Then
\[
 h_S(n)=\sum_{j=0}^2h_{\widehat K}(R^jn),
\]
and $S$ contains the full $R$-orbits of the four disks
\begin{equation}
 \mathbb D(A,2\eta),\quad\mathbb D(B,2\eta),
 \quad\mathbb D(C,2\eta),\quad\mathbb D(W,\eta),
 \qquad W=C+RA.                                  \label{eq:four-cap-disks}
\end{equation}
For a disk $\mathbb D(X,r)$ let
\[
 I(X,r)=\{n:\lVert n\rVert=1,\ \langle X,n\rangle+r\ge h\}.
\]
If the caps of the disks in \eqref{eq:four-cap-disks} cover the unit circle,
then $h_S(n)\ge h$ for every $n$ and \eqref{eq:lambda-support} proves
$\Lambda(\widehat K)\ge1$.

Put $\tau=h-2\eta=h(2c_*-1)$.  The rays of
\[
 A,\ B,\ C,\ W,\ RA                              \tag{6.1}
\]
occur in this cyclic order from $A$ to $RA$.  Indeed the two line distances
give
\[
 \frac{\operatorname{cross}(A,B)}{\lVert B-A\rVert}
 =\alpha(1-b)+\beta>0,
 \quad
 \frac{\operatorname{cross}(B,C)}{\lVert C-B\rVert}
 =\gamma+\delta(1-a)>0.
\]
If $A=(-X,-Y)$ and $C=(-P,-Q)$ in the centered cone basis, then
$0<X,Y,P,Q<1$, $X,Q>1/2$, and
\[
 \frac{\operatorname{cross}(C,RA)}h=PX+(Q-P)Y>0.
\]
The last sign is immediate unless $Q<P$ and $X<Y$, in which case it is
larger than $Q-P/2>0$.  The ray of $W=C+RA$ lies between those of $C$ and
$RA$.  The same coordinates give
\begin{equation}
 \lVert A\rVert,\lVert C\rVert>h/2>\eta.         \label{eq:outer-radii}
\end{equation}

For the two equal-radius overlaps it is enough to prove
\begin{equation}
 \frac{\operatorname{cross}(A,B)}{\lVert B-A\rVert}\ge\tau,
 \qquad
 \frac{\operatorname{cross}(B,C)}{\lVert C-B\rVert}\ge\tau.  \label{eq:adjacent-overlaps}
\end{equation}
We record the exact analytic verification.  By reflection take $a\le b$ and
set $m=1-b$, $M=1-a$, $U=a+b-1$, $s=m+M$, and $w=M-m$.  The smaller of the
two normalized left sides in \eqref{eq:adjacent-overlaps} is
\begin{equation}
 d=\frac{\mathcal A+U(2m-1)+D(\mathcal A+U)}{2\rho},
 \qquad \mathcal A=1-m+m^2.                      \label{eq:adjacent-d}
\end{equation}
On the $C_L$ sheet, write $F_L(z)=z^4-z^2+mz-m^2$.  This defining
polynomial evaluated at
$r=1-m/2+m^2/2$ is
\[
 F_L(r)=\frac{m^2(m-1)}{16}
 (m^5-3m^4+11m^3-17m^2+28m-12)>0.
\]
Indeed the polynomial in parentheses is increasing on $(0,1/2)$ and has
value $-33/32$ at $1/2$; its derivative is
$28-34m+m^2(5m^2-12m+33)>0$.  Since
$F_L(h)=-(m-h/2)^2\le0$ and the selected
root is unique, $2c_*-1\le\mathcal A$.  Put
\[
 X_0=\mathcal A+U,
 \qquad
 Y_0=2\rho\mathcal A-\mathcal A-U(2m-1).
\]
Then
\[
 d-\mathcal A=\frac{D(\mathcal A+U)-Y_0}{2\rho},
\]
and exact expansion gives
\[
\begin{aligned}
 Y_0={}&2(m^2-m+1)U^2+(2m^3-2m+3)U\\
      &+(m^2-m+1)(2m^2-2m+1)>0.
\end{aligned}
\]
Furthermore
\[
 D^2X_0^2-Y_0^2=4m\rho\Phi(U,m),
\]
where
\[
\begin{aligned}
 \Phi(U,m)={}&(1-m)(m^2-m+2)U^2\\
 &+(2-m^2)(m^2-m+1)U\\
 &-m(1-m)^2(m^2-m+1).
\end{aligned}
\]
In these variables
\[
 \chi=U^4-4U^3+5U^2-(m+2)U+m(1-m).
\]
Since $U^2(U^2-4U+5)>0$, the selector $\chi\le0$ implies
$U>m(1-m)/(m+2)$; $\Phi$ is increasing in $U$ and at that lower endpoint is
\[
 \frac{m^2(1-m)(m^4-2m^3+6m^2-6m+4)}{(m+2)^2}>0.
\]
The last quartic equals $(m^2-m)^2+5m^2-6m+4>0$.
Thus $d>\mathcal A\ge2c_*-1$.

On the other radial sheet let $E=\sqrt{4s^2-3}$.  Then
$0\le w<sE$ and $c_*=(s+w)/(1+E)$.  At fixed $U$, differentiation of
\eqref{eq:adjacent-d} can be checked without an implicit estimate.  Namely,
\[
 d=\frac{1+D}{2}-UB(D,w),\qquad
 B(D,w)=\frac{(1+U)(3+D)+w(1-D)}{D^2+3}.
\]
Here $D^2=w^2+6U+3U^2$, so $0\le D'=w/D\le1$, and
\[
 B_w=\frac{1-D}{D^2+3}\ge0.
\]
Writing $N=(1+U)(3+D)+w(1-D)$, direct differentiation gives
\[
 B_D=\frac{(1+U-w)(D^2+3)-2DN}{(D^2+3)^2}>-\frac{34}{27}.
\]
Indeed the first term in the numerator is positive, while
$1+U-w=2a>0$, $N<(7/6)4+1=17/3$, $D<1$, and $D^2+3\ge3$.  If $B_D<0$, then
$B_DD'\ge B_D$; otherwise $B_DD'\ge0$.  Hence
$B'=B_w+B_DD'>-34/27$, and $U<1/6$ yields
\[
 d'<\frac12+\frac16\frac{34}{27}=\frac{115}{162}<1,
\]
whereas
$(2c_*-1)'=2/(1+E)\ge1$.  Hence their difference decreases with $w$.
At the boundary $w=sE$, where $\chi=0$ and $c_*=s$, the preceding sheet
applies: this is an admissible point because $h<s<1$ and
$D^2=4s^4-12s+9<1$.  Indeed
$s^4-3s+2=(s-1)(s^3+s^2+s-2)<0$ on this interval; the second factor is
increasing and already equals $(7h-5)/4>0$ at $s=h$.  The preceding sheet gives
$d\ge2s-1$.  Therefore \eqref{eq:adjacent-overlaps} holds on both
sheets, and reflection supplies the other line.

It remains to check the two mixed overlaps.  On the two radial cells use the
rational upper envelopes
\[
 \bar c_L=s-\frac{\chi}{4s^3-2s+m},\qquad
 \bar c_T=s-\frac{\chi}{1-m}
\]
from Lemma~\ref{lem:global-branchwise-cbar}.  They satisfy
$c_*\le\bar c<1$, so
$\bar\eta=h(1-\bar c)$ obeys $0<\bar\eta\le\eta$.  The exact reduction in
Appendix~\ref{sec:exact-mixed-overlap} factors the two mixed Gram residuals
for the smaller caps into positive factors and eight integer-polynomial
signs.  Three fixed analytic charts and twenty global Bernstein identities
prove all eight signs without interval arithmetic or domain subdivision.
They are precisely the common-support inequalities for
\[
 \mathbb D(C,2\bar\eta),\ \mathbb D(W,\bar\eta)
 \quad\text{and}\quad
 \mathbb D(W,\bar\eta),\ \mathbb D(RA,2\bar\eta).
\]
Hence both pairs of smaller caps overlap.  Enlarging their radii from
$\bar\eta$ to $\eta$ proves the required two mixed overlaps.

Thus the four consecutive pairs of caps centered on the rays (6.1) overlap.
Every cap is a single arc of half-width less than $\pi/2$, and the positive
crosses put each overlap in the intervening short sector.  They cover the
sector from $A$ to $RA$; rotation through $2\pi/3$ covers the circle.  Hence
$h_S\ge h$, so $\Lambda(\widehat K)\ge1$.  Equation
\eqref{eq:tangent-reduction} completes the proof.
\end{proof}

\subsection{The center-independent obstruction}

\begin{theorem}[Center-independent direct nine-point obstruction]
\label{thm:zero-gap-obstruction}
There is no cover of $H$ by open unit equilateral triangles
$U_C,U_0,\ldots,U_5$ such that
\begin{enumerate}
\item every $U_i$ is a $\Vdzero$ role at $V_i$;
\item $U_0,\ldots,U_5$ cover $\partial H$; and
\item exactly one actual maximal boundary row is supercritical.
\end{enumerate}
No center class is assumed for $U_C$.
\end{theorem}

\begin{proof}
Lemma~\ref{lem:ab-extreme-jump} gives a parameter pair in
\eqref{eq:ab-strict-domain}.  Direct radial forcing and direct asymmetric
forcing give the compact containment $K_{\rm wit}(a,b)\subset U_C$.
Theorem~\ref{thm:witness-enclosure} says
$\Lambda(K_{\rm wit})\ge1$.

On the other hand, a compact subset of an open unit equilateral triangle has
positive clearance from its three sides.  Moving all three sides inward by a
sufficiently small equal amount gives a closed equilateral triangle of side
strictly less than one still containing the compact set.  Thus
$K_{\rm wit}\subset U_C$ would imply $\Lambda(K_{\rm wit})<1$, a
contradiction.
\end{proof}

\begin{remark}[Why exact-one is essential]
\label{rem:exact-one-essential}
The exact-one handoff theorem identifies the unique selected supercritical
row and makes every other selected row a strict descent.  Thus
\eqref{eq:extreme-chain} simultaneously provides the global minimum and
maximum used in the common radial demands, puts the two actual adjacent
handoffs in the intervals of Lemma~\ref{lem:fixed-line-signs}, and leaves only
one strict $AB$-frontier to analyze.  With two or more ascents this single
chain need not exist.  The $N_+\ge2$ branch is therefore handled by
Strategy~3 or Strategy~1.
\end{remark}

\begin{proposition}[Branches closed by direct nine-point forcing]
\label{prop:ab-core-branches}
Theorem~\ref{thm:zero-gap-obstruction} closes precisely the following
rows of the final routing table:
\begin{enumerate}
\item $\CEzero$, $N_+=1$, all six vertex roles $\Vdzero$;
\item $\CEone$ or $\CEtwo$, $N_+=1$, all six vertex roles $\Vdzero$, and
zero boundary gaps.
\end{enumerate}
\end{proposition}

\begin{proof}
In the first row, a missed boundary point would have to lie in the open
$\CEzero$ center role, giving a positive-length center trace on an incident
edge.  Hence the six vertex roles cover $\partial H$.  In the second row,
zero gaps means by definition that the two adjacent open vertex traces cover
every boundary edge, including every endpoint.  Thus in both rows all three
hypotheses of Theorem~\ref{thm:zero-gap-obstruction} hold.
\end{proof}
